%% COMP20008 Assignment 2 report, typeset in the elsarticle style of our arXiv paper.
%% Spec constraints kept: A4, single column, 11 pt, single line spacing, 2.54 cm margins.
\documentclass[final,1p,11pt]{elsarticle}
\geometry{a4paper,margin=2.54cm,footskip=1.2cm}
%% Times-like text and math without fontspec, so the file builds with pdfLaTeX.
\usepackage[T1]{fontenc}
\usepackage{newtxtext,newtxmath}
% inputenc leaves these three outside its default utf8 range
\DeclareUnicodeCharacter{2212}{\ensuremath{-}}
\DeclareUnicodeCharacter{2192}{\ensuremath{\rightarrow}}
\DeclareUnicodeCharacter{2264}{\ensuremath{\leq}}
\linespread{1.0}
\usepackage{graphicx}
\usepackage{booktabs,tabularx,array,makecell}
\usepackage[table]{xcolor}
\usepackage{microtype}
\usepackage{placeins}
\usepackage{xurl}
\usepackage[hidelinks]{hyperref}
\usepackage{titlesec}
\titleformat{\section}{\normalsize\bfseries}{\thesection.}{0.5em}{}
\titleformat{\subsection}{\normalsize\itshape}{\thesubsection.}{0.5em}{}
\titlespacing*{\section}{0pt}{8pt plus 2pt minus 2pt}{3pt}
\titlespacing*{\subsection}{0pt}{4pt plus 2pt minus 1pt}{1pt}

\definecolor{band}{HTML}{F3F5F8}
\definecolor{pick}{HTML}{DCE8F7}
\newcolumntype{L}{>{\raggedright\arraybackslash}X}
\renewcommand\theadfont{\bfseries}
\renewcommand\theadalign{bl}
\renewcommand{\arraystretch}{0.98}
\setlength{\tabcolsep}{4pt}
\setlength{\aboverulesep}{0.25ex}
\setlength{\belowrulesep}{0.35ex}
\setlength{\heavyrulewidth}{0.8pt}
\setlength{\abovecaptionskip}{2pt}
\setlength{\belowcaptionskip}{3pt}
\setlength{\textfloatsep}{6pt plus 2pt minus 2pt}
\setlength{\floatsep}{4pt plus 2pt minus 2pt}
\setlength{\intextsep}{4pt plus 2pt minus 2pt}
\setcounter{topnumber}{3}
\setcounter{bottomnumber}{2}
\setcounter{totalnumber}{5}
\renewcommand{\topfraction}{0.9}
\renewcommand{\bottomfraction}{0.6}
\renewcommand{\textfraction}{0.08}
\renewcommand{\floatpagefraction}{0.8}
\emergencystretch=3em
\renewcommand\_{\textunderscore\allowbreak}% long column names may break after an underscore
\tolerance=1200

%% plain page numbers instead of the "Preprint submitted to" footer
\makeatletter
\def\ps@pprintTitle{\let\@oddhead\@empty\let\@evenhead\@empty
  \def\@oddfoot{\hfil\footnotesize\thepage\hfil}\let\@evenfoot\@oddfoot}
\def\ps@plain{\let\@oddhead\@empty\let\@evenhead\@empty
  \def\@oddfoot{\hfil\footnotesize\thepage\hfil}\let\@evenfoot\@oddfoot}
\makeatother
\pagestyle{plain}

\begin{document}

\begin{center}
{\LARGE Property Features or Host Effort?}\\[4pt]
{\Large Inferring Superhost Status in the Melbourne Short-Term Rental Market}\\[10pt]
{\normalsize\textbf{Group W06G02} \quad COMP20008 Elements of Data Processing, Assignment 2}\\[3pt]
{\normalsize Jintao Liu (1582753)\quad Zhimo Ji (1810362)\quad Yifangfei Liu (1809764)\quad Leo Xu (1896121)}
\end{center}
\vspace{0pt}

\section{Introduction}

\noindent\textit{Research question.} Can superhost status be inferred from listing attributes and host operating behaviour once Airbnb's own criteria are excluded, and is it more closely associated with listing fundamentals or operational strategy?

Airbnb awards a 'superhost' status to hosts who meet a set of criteria. The criteria are at least 10 completed reservations (or 3 reservations totalling at least 100 nights), an overall rating of 4.8 or higher, a 90\% response rate to new messages within 24 hours, and a cancellation rate below 1\%, subject to specified exceptions (Airbnb, n.d.). The status is an indicator of exceptional service to guests and gives the listing more visibility, so hosts have a clear incentive to earn and keep it. 

We use the Inside Airbnb Melbourne dataset (Inside Airbnb, 2026), containing 25,728 listings described by 90 columns. The days\_since\_last\_review reference date is 28/06/2026, which is also the latest last\_review in the data.

\section{Methodology}

\subsection{Data Preprocessing}

We proposed and considered six candidate preprocessing steps (Table 1) for the Melbourne Airbnb dataset, aligned with the research question.

\begin{table}[!htbp]
\centering
\caption{Candidate preprocessing steps with dataset properties that supported the selection condition, with the corresponding outcome.}
\footnotesize
\begin{tabularx}{\textwidth}{p{0.25cm}>{\raggedright\arraybackslash}p{4.9cm}Lp{1.3cm}}
\toprule
\textbf{\#} & \textbf{Candidate step} & \textbf{Properties} & \textbf{Decision} \\
\midrule
1 & Deriving days\_since\_last\_review from the raw last\_review field & 4,477 rows (17.4\%) have no review date at all, and their superhost rate is 6.59\% & Selected \\
\rowcolor{band}2 & Counting amenity\_count by parsing the structured amenities field & Superhosts average 44.08 amenities, non-superhosts 31.03 & Selected \\
3 & Imputing missing price values via median imputation plus a missing indicator variable & 6,553 rows (25.47\%) have no price & Selected \\
\rowcolor{band}4 & Imputing missing values in the bedrooms and beds fields & Missing 18.19\% for bedrooms and 27.44\% for beds, and the missingness is conditional on room type & Rejected \\
5 & Merging rare categories for the property\_type variable & Of the 82 categories, 41 contain fewer than 10 listings, and these 41 together account for only 111 rows (0.43\%) & Rejected \\
\rowcolor{band}6 & Splitting minimum\_nights into short/long groups at a 30-night threshold & Only 426 rows (1.66\%) require 30 nights or more & Rejected \\
\bottomrule
\end{tabularx}
\end{table}

\textit{Justifications for selections.} Deriving days\_since\_last\_review from the raw last\_review field was selected as dropping the 4,477 listings with no review date (17.4\% of all rows) would raise the superhost rate from 30.65\% to 35.72\%. Furthermore, 295 superhost listings are still in the no review date group, suggesting the rows are worth keeping.

Amenities are also among the few features a host can change at any time, which places this column on the operational side of the research question.

Handling missing price values via imputation was selected as price missingness depends on whether the listing is bookable at all and is missing at random (MAR). 99.85\% of unavailable listings (availability\_365 = 0) have no price, against 2.22\% of available ones. As unavailable listings have an 8.13\% superhost rate against 37.68\% for available listings, deletion would inflate the superhost rate from 30.65\% to 38.06\%.

\textit{Justifications for rejections.} Imputing the bedrooms and beds fields was rejected because missingness is structural by room type. Bedrooms are missing for 18.19\% (beds 27.44\%), but the data is split by room type, with entire homes missing 2.47\%, private rooms 61.28\%. Thus, we observe that the blanks are structural. In a private room the bedroom count is usually absent, and in a whole home it almost never is.

Merging property\_type categories was rejected because although there are 41 property types having fewer than 10 listings (out of 82 total property types), these 41 account for only 111 rows (0.43\%), which is a trivial change to the data (0.43\% of 25,728). The property\_type field is excluded from the models, as we will derive a new field, room\_type\_3, that is detailed within model training.

Splitting minimum\_nights into short/long stay bins was rejected as only 426 listings (1.66\%) require a minimum stay of 30 nights or more. The long-stay group would be too small for the model to learn from.

\subsection{Correlation Analysis}

\textit{Variable set.} The research question asks whether superhost status is more closely associated with listing fundamentals or host operational strategy, so the exploratory variable set must include representatives on both sides.

The five variables selected are host\_listings\_count, days\_since\_last\_review, availability\_365, dist\_cbd, and the binary target y ($y = 1$ when host\_is\_superhost = 't', $y = 0$ otherwise). Five variables give 10 unique pairs and 40 values.

\textit{Rationale behind variable set.} On the listing fundamentals side, dist\_cbd measures the distance from Melbourne CBD (Flinders Street Station) and is appropriate as a fixed physical attribute that no host can change.

Price is excluded because 25.47\% of prices are missing and 93.35\% of those missing rows are unavailable listings, so keeping only priced rows would give a non-random sample. On the fundamentals side, the only other remaining feature, accommodates, has the lowest mutual information among the 11 features (0.0093), so dist\_cbd is used as the single fundamentals representative.

On the operational strategy side, the three variables selected characterise hosts based on three dimensions. host\_listings\_count captures scale (median 2, maximum 667), availability\_365 reflects calendar positioning (21.8\% of listings at 0 days, 1.56\% at 365) and days\_since\_last\_review measures recent activity (median 84 days, maximum 4,593). The variables are biased 3:1 towards operational features, as only dist\_cbd represents fundamentals, so the fair comparison between the two families relies on feature-subset experiments rather than this variable set.

\textit{Computation.} The four exploratory variables are continuous and the target is binary, but no single coefficient is equally suitable for every pair. Applied to a 0/1 target, Pearson is a point-biserial correlation and reflects only the difference in means between the two groups, rather than the shape of the relationship (Tate, 1954). Spearman leaves the target with just two ranks and many ties. For MI and NMI the continuous variables are binned using equal-frequency bins. Because MI rises with the bin count, so the count cannot be chosen by maximising it. Across 5, 10 and 20 bins the accompanying NMI for days\_since\_last\_review $\sim$ availability\_365 varies from 0.0799 to 0.0856 to 0.0896, so we chose $n=10$ as the midrange value. Equal-width binning was rejected as 10 equal-width bins would concentrate 91.4\% of host\_listings\_count rows in one bin (84.6\% at 20 bins), which lowers MI. Correlations were only computed on 21,251 listings with at least one review, because 4,477 zero-review listings carry the imputation fill value 4,594 for days\_since\_last\_review, which would distort the coefficients because they all share the same artificial maximum value.

\subsection{Supervised Learning Models and Evaluation}

\textit{Data.} All 25,728 listings are used, with zero-review listings marked by has\_review. Every listing held by the same host carries the same label. 3,770 of 14,113 hosts are superhosts (26.71\%), but because superhosts hold more listings each they account for 30.65\% of rows. The data is therefore split by host group.

\textit{Target and features.} The 11 features comprise three variables related to fundamentals (room\_type\_3, dist\_cbd, accommodates) and eight related to operational features (availability\_365, amenity\_count, days\_since\_last\_review, has\_review, price\_num, price\_missing, minimum\_nights, host\_listings\_count). room\_type\_3 merges the 231 shared-room and 52 hotel-room rows (283 in total, 1.1\%) into Other, leaving three categories that are then one-hot encoded. price\_num and price\_missing are treated as operational because the host sets the price, although price\_num and accommodates (a listing fundamental) have a Spearman correlation of 0.682. We justify this by moving price into the fundamentals group, giving a CV F1 score of 0.1844 for the decision tree and 0.2647 for k-NN, against 0.5577 and 0.5442 for the operational group. Thus the conclusion is unchanged.

\textit{Test/train split.} We use StratifiedGroupKFold, grouping by host\_id and stratifying by y, and take one fold as the test set, giving an 80/20 split. The 20,582 training rows contain 6,308 superhost listings and 14,274 non-superhost listings from 11,292 hosts, while the 5,146 test rows contain 1,577 and 3,569 listings from 2,821 hosts, with no host on both sides. The superhost share is held at 30.65\% on each side. Grouping is necessary because the label belongs to the host and 55.8\% of listings share a host with at least one other listing. Under a random split, a model could score well by recognising a host it had already seen, inflating performance metrics.

\textit{Validation strategy.} Hyperparameters are tuned on the 20,582 training rows using 5-fold, host-grouped and stratified cross-validation, and each fold is scored by F1 for the superhost class. Each validation fold holds about 4,116-4,117 rows, including 1,261-1,262 superhost listings from 2,258-2,259 hosts. Ten folds were not used because the largest host holds 319 listings. That host would occupy 7.75\% of a validation fold under five folds and 15.5\% under ten, which means a single host could dominate a whole fold.

\textit{Hyperparameter grids.} k-NN is trained with k $\in$ \{1, 5, 11, 21, 51, 101, 151, 201, 301\} and the decision tree with max\_depth $\in$ \{2, 3, 4, 5, 6, 8, 10, 15, None\}, where None means unlimited depth and is equivalent to memorising the training set. Only odd values of k are tested to avoid ties during classification.

\textit{Baseline and metrics.} The primary metric is the F1 score of the superhost class, and per-class precision and recall and overall accuracy are reported alongside. The baseline is a majority-class classifier, which predicts 'not a superhost' for every listing. It reaches 0.694 accuracy and an F1 score of exactly 0. We perform bootstrapping with 1,000 resamples drawn with replacement that  gives a 95\% confidence interval for test F1.  

\textit{Feature-set experiments.} To answer the RQ, both models are retrained on six labelled feature sets. A uses the three fundamentals features, B the eight operational features, C all eleven, D removes the review pair from C, E removes availability\_365 from C, and F removes the review pair from B. Two further checks, G (C without amenity\_count) and H (C without the price pair), support the preprocessing discussion. Each set uses the hyperparameters selected in Table 5 and is scored by cross-validation on the training rows, with the test score used only for confirmation.

\subsection{Feature Selection}

\textit{Filter method.} Mutual information with the target is estimated independently for each feature on the training set using scikit-learn's mutual\_info\_classif function. Room type is integer encoded, the discrete features (room\_type\_3, has\_review, price\_missing) are declared, and remaining missing values are filled with the training median. These continuous-value MIs are not comparable to the binned MI in Table 4.

\textit{Embedded method.} Decision-tree impurity is taken from the tuned depth-10 decision tree and the tree is trained on the training set only.

\subsection{Dimensionality Reduction \& Clustering}

\textit{Candidate feature sets.} Clustering runs on the 21,251 reviewed listings, since the 4,477 zero-review listings share the imputed value 4,594, producing a separate cluster. Three candidate feature sets are compared in Table 2: (S1) operational only, (S2) fundamentals only, and (S3) all six features. Three of the six features are log-transformed via log1p, and all six were standardised via StandardScaler. Without log transform a single host with 667 listings would dominate a component on its own.

\begin{table}[!htbp]
\centering
\caption{Candidate feature sets for clustering, built from the two families in the research question. Every candidate set follows the same pipeline of log1p transformation, standardisation and K-Means with $k = 5$. The superhost-rate range is the difference between the highest and lowest cluster superhost rates, and the smallest cluster is the number of listings in the least populated cluster.}
\footnotesize
\begin{tabularx}{\textwidth}{p{0.5cm}>{\raggedright\arraybackslash}p{2.3cm}L>{\raggedleft\arraybackslash}p{1.8cm}>{\raggedleft\arraybackslash}p{1.45cm}}
\toprule
\textbf{Set} & \textbf{Name} & \textbf{Features} & \thead{Superhost-\\rate range} & \thead{Smallest\\ cluster} \\
\midrule
S1 & Operational only & availability\_365, amenity\_count, days\_since\_last\_review, host\_listings\_count & 0.5331 & 3,243 \\
\rowcolor{band}S2 & Fundamentals only & dist\_cbd, accommodates & 0.0754 & 780 \\
S3 & Mixed (both families) & availability\_365, amenity\_count, days\_since\_last\_review, host\_listings\_count, dist\_cbd, accommodates & 0.4719 & 2,365 \\
\bottomrule
\end{tabularx}
\end{table}

\textit{Selected feature set.} Feature set S3 is used because the research question asks whether superhost status is closer to listing fundamentals or to operational strategy, and a set holding only one family cannot answer that. S1 separates slightly better but contains operational features alone, so it cannot test which family matters, whereas S3 lets both families compete while retaining most of S1's separation (a superhost-rate range of 0.4719 versus 0.5331). The very small range for S2 (0.0754) indicates that clustering on fundamentals alone produces groups that barely differ in superhost status.

\textit{Clustering and PCA.} $k=5$ is read from the bend in the elbow curve (Figure 1). The SSE reduction when adding each cluster is 16.0\%, 12.8\%, 12.0\% and 6.8\% for k=2 to 3, 3 to 4, 4 to 5 and 5 to 6. For hierarchical clustering, Ward linkage with Euclidean distance is used at the same $k=5$. 

We justify Ward linkage as single, average, and complete linkages all collapse these features at $k=5$ into one dominant cluster. Single linkage gives sizes [21247, 1, 1, 1, 1], average linkage [19482, 1401, 341, 26, 1], and complete linkage a largest cluster of 12,473. Ward linkage returns the most balanced partition (largest cluster 8,256) and is thus chosen. PCA is applied to feature set S3 with all six components, and the explained variance of the first three is reported.

\begin{figure}[!htbp]
\centering
\includegraphics[width=0.5\textwidth]{outputs/fig_elbow.png}
\caption{K-Means SSE by number of clusters $k$ on feature set S3 (n = 21,251 reviewed listings). The dashed line marks the selected $k = 5$.}
\end{figure}

\section{Results Exploration and Analysis}

\subsection{Data Preprocessing}

\begin{table}[!htbp]
\centering
\caption{Measurable before/after impact of the three selected preprocessing steps.}
\footnotesize
\begin{tabularx}{\textwidth}{>{\raggedright\arraybackslash}p{3.4cm}LL}
\toprule
\textbf{Candidate step} & \textbf{Before} & \textbf{After} \\
\midrule
Deriving days\_since\_last\_review from last\_review & last\_review missing for 4,477 rows, could not derive missing values & days\_since\_last\_review filled with 4,594 days for missing last\_review rows. 0 missing rows with median 84 days. \\
\rowcolor{band}Counting amenity\_count by parsing the structured amenities field & No field exists describing the amenity count of a particular listing. & A naive comma count would be wrong on 3,096 rows (12.03\%), over-counting by up to 21 amenities for a single listing. \\
Imputing missing price values via median imputation plus a missing indicator variable & 25.47\% of prices missing; the missingness depends on availability\_365 (MAR, not MCAR) & 0\% prices missing, imputed within the pipeline using the training-set median of \$245.00. price\_missing marks every imputed row. \\
\bottomrule
\end{tabularx}
\end{table}

\subsection{Correlation Analysis}

\begin{table}[!htbp]
\centering
\caption{Pearson, Spearman, mutual information (bits) and normalised mutual information for each pair of the five variables (n=21,251 listings with at least one review). Rows are ordered by NMI, highest first.}
\footnotesize
\begin{tabularx}{\textwidth}{L>{\raggedleft\arraybackslash}p{1.4cm}>{\raggedleft\arraybackslash}p{1.4cm}>{\raggedleft\arraybackslash}p{1.4cm}>{\raggedleft\arraybackslash}p{1.4cm}}
\toprule
\textbf{Pair} & \textbf{Pearson} & \textbf{Spearman} & \textbf{MI (bits)} & \textbf{NMI} \\
\midrule
days\_since\_last\_review \textasciitilde{} availability\_365 & −0.438 & −0.395 & 0.2639 & 0.0856 \\
\rowcolor{band}days\_since\_last\_review \textasciitilde{} y & −0.337 & −0.443 & 0.1672 & 0.0785 \\
availability\_365 \textasciitilde{} y & 0.146 & 0.155 & 0.0613 & 0.0324 \\
\rowcolor{band}host\_listings\_count \textasciitilde{} dist\_cbd & −0.131 & −0.225 & 0.0837 & 0.0285 \\
host\_listings\_count \textasciitilde{} days\_since\_last\_review & −0.113 & −0.194 & 0.0806 & 0.0274 \\
\rowcolor{band}host\_listings\_count \textasciitilde{} availability\_365 & 0.059 & 0.225 & 0.0621 & 0.0230 \\
availability\_365 \textasciitilde{} dist\_cbd & 0.170 & 0.124 & 0.0455 & 0.0147 \\
\rowcolor{band}days\_since\_last\_review \textasciitilde{} dist\_cbd & −0.081 & 0.061 & 0.0458 & 0.0138 \\
host\_listings\_count \textasciitilde{} y & −0.069 & 0.070 & 0.0169 & 0.0097 \\
\rowcolor{band}dist\_cbd \textasciitilde{} y & 0.076 & 0.009 & 0.0101 & 0.0047 \\
\bottomrule
\end{tabularx}
\end{table}

Pearson and Spearman disagree on host\_listings\_count with the target, at −0.069 and +0.070, because the superhost rate does not move in one direction. This pattern is shown in Figure 2.

\begin{figure}[!htbp]
\centering
\includegraphics[width=0.46\textwidth]{outputs/fig_hlc_vs_y.png}
\caption{Superhost rate by host\_listings\_count band (n = 21,251 reviewed listings). The red line marks the overall superhost rate of 35.7\%.}
\end{figure}

\subsection{Supervised Learning Models and Evaluation}

\begin{table}[!htbp]
\centering
\caption{Cross-validated F1 score at every hyperparameter value tried (five grouped folds, $n=20,582$). Selected configurations shaded blue. SD denotes standard deviation.}
\footnotesize
\begin{tabularx}{\textwidth}{>{\raggedright\arraybackslash}p{2.6cm}*{9}{>{\centering\arraybackslash}X}}
\toprule
\textbf{k-NN: k} & 1 & 5 & 11 & 21 & 51 & \cellcolor{pick}\textbf{101} & 151 & 201 & 301 \\
CV F1 (mean) & 0.4802 & 0.5020 & 0.5071 & 0.5190 & 0.5299 & \cellcolor{pick}\textbf{0.5370} & 0.5312 & 0.5309 & 0.5249 \\
\rowcolor{band}$\pm$ SD & 0.0175 & 0.0159 & 0.0176 & 0.0112 & 0.0121 & 0.0133 & 0.0147 & 0.0175 & 0.0226 \\
\midrule
\textbf{Tree: max\_depth} & 2 & 3 & 4 & 5 & 6 & 8 & \cellcolor{pick}\textbf{10} & 15 & None \\
CV F1 (mean) & 0.5158 & 0.5611 & 0.5712 & 0.5384 & 0.5512 & 0.5622 & \cellcolor{pick}\textbf{0.5745} & 0.5363 & 0.5018 \\
\rowcolor{band}$\pm$ SD & 0.0503 & 0.0431 & 0.0177 & 0.0424 & 0.0240 & 0.0123 & 0.0130 & 0.0171 & 0.0072 \\
\bottomrule
\end{tabularx}
\end{table}

\begin{table}[!htbp]
\centering
\caption{Change in mean 5-fold CV F1 when each tuned hyperparameter is moved from its default to the selected value.}
\footnotesize
\begin{tabularx}{\textwidth}{>{\raggedright\arraybackslash}p{1.9cm}L>{\raggedright\arraybackslash}p{2.3cm}>{\raggedright\arraybackslash}p{2.3cm}>{\raggedleft\arraybackslash}p{1.3cm}>{\raggedleft\arraybackslash}p{1.3cm}>{\raggedleft\arraybackslash}p{1.3cm}}
\toprule
\textbf{Model} & \textbf{Hyperparameter} & \textbf{Default} & \textbf{Chosen} & \thead{Default\\ CV F1} & \thead{Chosen\\ CV F1} & \textbf{Effect} \\
\midrule
k-NN & Number of neighbours (k) & k = 5 & k = 101 & 0.5020 & 0.5370 & +0.0350 \\
\rowcolor{band}Decision Tree & Maximum depth of tree (max\_depth) & max\_depth = None & max\_depth 
= 10 & 0.5018 & 0.5745 & +0.0727 \\
\bottomrule
\end{tabularx}
\end{table}

\begin{table}[!htbp]
\centering
\caption{Test-set performance (n=5,146), with precision, recall, accuracy and F1 scores.}
\footnotesize
\begin{tabularx}{\textwidth}{>{\raggedright\arraybackslash}p{3.8cm}L>{\raggedleft\arraybackslash}p{1.7cm}>{\raggedleft\arraybackslash}p{1.7cm}>{\raggedleft\arraybackslash}p{1.7cm}>{\raggedleft\arraybackslash}p{1.7cm}}
\toprule
\textbf{Model} & \textbf{Class} & \textbf{Precision} & \textbf{Recall} & \textbf{Accuracy} & \textbf{F1 Score} \\
\midrule
Baseline & Non-superhost & 0.6935 & 1.0000 & 0.6935 & 0.8190 \\
 & Superhost & 0.0000 & 0.0000 &  & 0.0000 \\
k-NN (k=101) & Non-superhost & 0.8106 & 0.8322 & 0.7487 & 0.8212 \\
 & Superhost & 0.5958 & 0.5599 &  & 0.5773 \\
Decision Tree (depth=10) & Non-superhost & 0.8445 & 0.8067 & 0.7629 & 0.8252 \\
 & Superhost & 0.6028 & 0.6639 &  & 0.6319 \\
\bottomrule
\end{tabularx}
\end{table}

\begin{table}[!htbp]
\centering
\caption{95\% percentile-bootstrap intervals (1,000 paired host resamples) for test-set superhost-class F1, for the gap between models and for feature-subset comparisons. A, B and F are decision-tree results.}
\footnotesize
\begin{tabularx}{\textwidth}{L>{\raggedleft\arraybackslash}p{3.6cm}>{\raggedleft\arraybackslash}p{2.6cm}}
\toprule
\textbf{Model/Pair} & \textbf{95\% CI} & \textbf{Point estimate} \\
\midrule
k-NN & [0.5423, 0.6155] & 0.5773 \\
\rowcolor{band}Decision Tree & [0.5940, 0.6757] & 0.6319 \\
DT − k-NN & [0.0130, 0.0986] & 0.0546 \\
\rowcolor{band}A (fundamentals) & [0.1535, 0.2384] & 0.1955 \\
B (operational) & [0.5981, 0.6844] & 0.6389 \\
\rowcolor{band}B − A & [0.3813, 0.5046] & 0.4434 \\
F (operational, no review pair) & [0.4923, 0.6313] & 0.5662 \\
\rowcolor{band}F − A & [0.2819, 0.4536] & 0.3707 \\
\bottomrule
\end{tabularx}
\end{table}

\FloatBarrier

\subsection{Feature Selection}

\textit{Top features from each method.} Table 9 reports the full feature ranking under mutual information and under tree impurity, with the corresponding values.

\textit{Hard case.} Listing 6812677, which is from host 35667203, is one of the tree's 530 false negatives and among its most confident. The superhost listing was scored 0.000 by the tree, the class share of a leaf holding 86 training listings with no superhosts.

\begin{table}[!htbp]
\centering
\caption{Full ranking of the 11 features by mutual information and decision-tree impurity, with permutation importance for both models.}
\footnotesize
\begin{tabularx}{\textwidth}{L>{\raggedleft\arraybackslash}p{1.3cm}>{\raggedleft\arraybackslash}p{1.0cm}>{\raggedleft\arraybackslash}p{1.5cm}>{\raggedleft\arraybackslash}p{1.4cm}>{\raggedleft\arraybackslash}p{2.2cm}>{\raggedleft\arraybackslash}p{2.2cm}}
\toprule
\textbf{Feature} & \thead{MI\\(bits)} & \thead{MI\\ rank} & \textbf{Impurity} & \thead{Impurity\\ rank} & \thead{DT\\ permutation} & \thead{k-NN\\ permutation} \\
\midrule
days\_since\_last\_review & 0.1837 & 1 & 0.5005 & 1 & 0.2269 & 0.0747 \\
\rowcolor{band}host\_listings\_count & 0.1569 & 2 & 0.2698 & 2 & 0.0195 & 0.0037 \\
availability\_365 & 0.1095 & 3 & 0.0356 & 4 & 0.0131 & 0.0169 \\
\rowcolor{band}amenity\_count & 0.1080 & 4 & 0.1023 & 3 & 0.0810 & 0.1433 \\
price\_num & 0.0948 & 5 & 0.0277 & 6 & 0.0109 & 0.0039 \\
\rowcolor{band}price\_missing & 0.0651 & 6 & 0.0024 & 9 & 0.0219 & 0.0621 \\
has\_review & 0.0496 & 7 & 0.0000 & 11 & 0.0000 & 0.0073 \\
\rowcolor{band}dist\_cbd & 0.0490 & 8 & 0.0278 & 5 & 0.0054 & 0.0324 \\
minimum\_nights & 0.0140 & 9 & 0.0157 & 8 & 0.0014 & 0.0006 \\
\rowcolor{band}room\_type\_3 & 0.0128 & 10 & 0.0001 & 10 & 0.0000 & 0.0168 \\
accommodates & 0.0093 & 11 & 0.0181 & 7 & 0.0062 & -0.0064 \\
\bottomrule
\end{tabularx}
\end{table}

\subsection{Dimensionality Reduction \& Clustering}

\textit{PCA.} The first three components carry 69.4\% of the variance between them, and Table 10 reports the explained variance and the top loadings. Figure 3 plots the listings against the first two components.

\begin{table}[!htbp]
\centering
\caption{PCA on standardised set S3, containing explained variance ratio (EVR), cumulative variance and top 3 loadings.}
\footnotesize
\begin{tabularx}{\textwidth}{p{1.6cm}>{\raggedleft\arraybackslash}p{0.95cm}>{\raggedleft\arraybackslash}p{1.6cm}>{\raggedright\arraybackslash}p{4.0cm}LL}
\toprule
\textbf{Component} & \textbf{EVR} & \textbf{Cumulative} & \textbf{Rank 1} & \textbf{Rank 2} & \textbf{Rank 3} \\
\midrule
PC1 & 0.3211 & 0.3211 & days\_since\_last\_review +0.563 & amenity\_count −0.504 & availability\_365 −0.447 \\
\rowcolor{band}PC2 & 0.2045 & 0.5255 & dist\_cbd +0.786 & host\_listings\_count −0.556 & accommodates +0.208 \\
PC3 & 0.1681 & 0.6936 & accommodates +0.671 & availability\_365 −0.591 & amenity\_count +0.355 \\
\bottomrule
\end{tabularx}
\end{table}

\begin{figure}[!htbp]
\centering
\includegraphics[width=0.47\textwidth]{outputs/fig_pca_clusters.png}
\caption{PCA scatter of the 21,251 reviewed listings on the first two components, coloured by superhost status, with K-Means centroids (k=5) marked. A random sample of 5,000 listings is shown. Centroids numbered 1-5 correspond to clusters K1-K5 in Table 11.}
\end{figure}

\begin{table}[!htbp]
\centering
\caption{K-Means (K1-K5) and agglomerative clustering using Ward (W1-W5) on the standardised six-feature set with $k=5$. Feature values are medians, and are not standardised.}
\footnotesize
{\setlength\tabcolsep{2.6pt}\begin{tabular}{>{\raggedright\arraybackslash}p{14pt}>{\raggedleft\arraybackslash}p{36pt}>{\raggedleft\arraybackslash}p{30pt}>{\raggedleft\arraybackslash}p{44pt}>{\raggedleft\arraybackslash}p{40pt}>{\raggedleft\arraybackslash}p{33pt}>{\raggedleft\arraybackslash}p{54pt}>{\raggedleft\arraybackslash}p{38pt}>{\raggedleft\arraybackslash}p{48pt}>{\raggedleft\arraybackslash}p{46pt}}
\toprule
\thead{\#} & \thead{Listings} & \thead{Share} & \thead{Superhost\\rate} & \thead{dist\_cbd} & \thead{accom-\\modates} & \thead{availability\\\_365} & \thead{amenity\\\_count} & \thead{days\_since\\\_last\\\_review} & \thead{host\\\_listings\\\_count} \\
\midrule
\textbf{K1} & 4,541 & 21.4\% & 35.3\% & 1.42 & 4 & 188 & 41 & 50 & 50 \\
\rowcolor{band}\textbf{K2} & 5,340 & 25.1\% & 53.9\% & 2.25 & 4 & 157 & 47 & 40 & 2 \\
\textbf{K3} & 4,689 & 22.1\% & 6.8\% & 4.54 & 2 & 0 & 23 & 1,883 & 1 \\
\rowcolor{band}\textbf{K4} & 2,365 & 11.1\% & 43.7\% & 20.56 & 8 & 218 & 51 & 58 & 3 \\
\textbf{K5} & 4,316 & 20.3\% & 40.7\% & 17.82 & 2 & 310 & 33 & 76 & 2 \\
\rowcolor{band}\textbf{W1} & 8,256 & 38.9\% & 49.9\% & 5.36 & 4 & 174 & 47 & 56 & 1 \\
\textbf{W2} & 3,829 & 18.0\% & 30.1\% & 7.32 & 2 & 267 & 22 & 100 & 3 \\
\rowcolor{band}\textbf{W3} & 4,716 & 22.2\% & 37.5\% & 1.79 & 4 & 175 & 42 & 55 & 45 \\
\textbf{W4} & 3,249 & 15.3\% & 2.9\% & 4.34 & 2 & 0 & 20 & 2,364 & 1 \\
\rowcolor{band}\textbf{W5} & 1,201 & 5.7\% & 37.8\% & 17.55 & 10 & 221 & 49 & 49 & 12 \\
\bottomrule
\end{tabular}}
\end{table}

\section{Discussion and Interpretation}

\subsection{Data Preprocessing}

The three preprocessing steps have different effects on the models, as shown by the feature-subset experiments. Removing the review-related features (experiment D) moves the decision tree's CV F1 from 0.5745 to 0.5262 (test F1 0.6319 to 0.5667) and k-NN from 0.5370 to 0.5032. Removing amenity\_count (experiment G) moves k-NN from 0.5370 to 0.4747 (test 0.5773 to 0.5196), while the decision tree is almost unchanged at 0.5747. Removing the price pair (experiment H) moves k-NN to 0.5220 and the decision tree slightly up to 0.5791. We conclude that the decision tree heavily relies on review recency, whereas k-NN also depends on amenity\_count and price, which matches the permutation-importance results in Table 9.

\subsection{Correlation Analysis}

\textit{Agreement and divergence.} The two coefficients agree closely on the strongest pair, days\_since\_last\_review with availability\_365 (−0.438 vs −0.395), but they point in opposite directions on two pairs (Table 4).

The host\_listings\_count result can be explained from the data. Spearman works with ranks, with most listings fall in the 0-5 host-listings range (8,103, 2,560 and 2,832 listings across its three bands), where the superhost rate rises from 28.5\% to 46.1\%, so the rank correlation is positive at +0.070. Pearson uses the value instead, with 1,259 listings belong to hosts with more than 100 listings, and their superhost rate is only 16.9\%, yielding a coefficient of -0.069. MI and NMI do not assume a particular shape, yet NMI for this pair is only 0.0097.

\textit{Sign reversal.} Between days\_since\_last\_review and dist\_cbd the same pattern appears (Pearson -0.081, Spearman +0.061). When distance is split into twenty equal-frequency bands, the median days\_since\_last\_review is 41-60 days at 0.6-1.4 km, rises to 111 days at 1.7 km, peaks at 122 days at 5.0 km, and then falls to 69 days at 34.6 km. Across the inner and middle bands, days\_since\_last\_review rises with distance, which is what the positive Spearman value measures. The outer bands are both far from the centre, causing the Pearson value to be negative.

\textit{Connection to the RQ.} Of the four pairs ranked by NMI, days\_since\_last\_review has the strongest association (NMI 0.0785), followed by availability\_365 (0.0324) and host\_listings\_count (0.0097). dist\_cbd comes last (0.0047), so the only listing fundamentals feature in the set has the weakest association with the target. It is important to note that these are associations only. Hosts with more than 100 listings may be running business or apartment models, so their low rate may reflect the business, rather than service quality. Not every guest also leaves a review, so days\_since\_last\_review may be inaccurate. A further confounder is that days\_since\_last\_review partly substitutes for an excluded criterion, since its Spearman correlation with number\_of\_reviews\_ltm is -0.8155. This motivates experiment F, which removes the review-related operational features. The test-set interval for F - A, [0.2819, 0.4536], remains positive (Table 8), meaning operational features still predict better than fundamentals.

\subsection{Supervised Learning Models and Evaluation}

\textit{Model performance.} With the official criteria excluded, the tree reaches a test F1 of 0.632 on the superhost class and k-NN 0.577, against a majority-class baseline of exactly 0. The baseline model that ignores every feature already reaches 0.694 accuracy, which is just the non-superhost share, so accuracy cannot separate a useful model from the baseline. Relative to the baseline, F1 rises by 0.6319 for the tree and 0.5773 for k-NN, accuracy by 0.0694 and 0.0552 (Table 7).

\textit{Comparing the models.} The decision tree is the better model of the two. Its paired bootstrap interval against k-NN is [0.0130, 0.0986], which does not include zero, so the difference is unlikely to be sampling noise. The tree predicts more listings as superhosts, 1,737 in total against 1,482 for k-NN, which is a disadvantage (1,047 and 883 are true positives, 690 and 599 are false positives).

\textit{Feature influence.} The two models rely on different features. For the decision tree, permutation importance is highest for days\_since\_last\_review (0.2269), followed by amenity\_count (0.0810). For k-NN it is highest for amenity\_count (0.1433), followed by days\_since\_last\_review (0.0747) and price\_missing (0.0621). host\_listings\_count shows a contrast, with impurity importance of 0.2698 (second highest) but permutation importance of only 0.0195 for the tree and 0.0037 for k-NN, meaning the tree often splits on it during training although it adds little to test predictions. The features that both models depend on most are operational, which is consistent with the RQ.

\textit{Model behaviour and limits.} Depth 10 and depth 4 differ by 0.0033 in CV F1 (0.5745 $\pm$ 0.0130 against 0.5712 $\pm$ 0.0177), less than either standard deviation, so the two cannot be separated reliably. Depth 10 is chosen as it has higher mean CV F1. The decision tree also overfits mildly, with a F1 score of 0.7188 on training against 0.6319 on test (a gap of 0.087). k-NN narrows that to 0.6023 against 0.5773.

\textit{Subset experiments and the RQ.} The subset experiments answer the research question most directly. On cross-validation, fundamentals reach CV F1 scores of 0.166 and operational features 0.548. On the test set the corresponding values are 0.1955 and 0.6389, and the 95\% interval for the test-set difference of 0.4434 is [0.3813, 0.5046], which does not include zero.

\subsection{Feature Selection}

\textit{Agreements and divergence.} Both methods (MI and impurity) put days\_since\_last\_review first and host\_listings\_count second, but they swap the next two features: availability\_365 ranks third by mutual information (0.1095) but fourth by impurity (0.0356), while amenity\_count ranks fourth by mutual information (0.1080) but third by impurity (0.1023). Mutual information scores each feature against the target on its own, so the redundancy between availability\_365 and days\_since\_last\_review is not penalised. Impurity works differently, measuring a feature by the impurity it removes at tree splits. availability\_365 carries much the same information as days\_since\_last\_review, so once the tree has split on recency there is little left for availability\_365 to add, and it drops to 0.0356. amenity\_count overlaps with recency far less, so it keeps 0.1023. Experiment E supports this reading. Removing availability\_365 leaves CV F1 slightly higher (0.5846 vs 0.5745) and test F1 essentially unchanged (0.6303 vs 0.6319).

\textit{The hard case.} Listing 6812677 is a superhost listing that the tree classed as a non-superhost with a predicted probability of 0.000. On the features that trained the model it looks less active: 18 amenities against medians of 45 for superhosts and 32 for non-superhosts, availability of 116 which is below even the non-superhost median of 124, a host with a single listing, and 156 days since the last review against a superhost median of 37. The tree path starts with the root split days\_since\_last\_review $\leq$ 147.5, then meets host\_listings\_count $\leq$ 1.5 and amenity\_count $\leq$ 24.5, ending in a leaf of 86 training listings that are all non-superhosts, yielding 0.000. k-NN also classifies it incorrectly, with a probability of 0.267. On the two official criteria observable in the data, it meets both: 24 reviews in the last twelve months (at least 10 required) and a rating of 4.8 (at least 4.8 required). Its days\_since\_last\_review feature is at the 90.8th percentile among reviewed superhosts, so it is extreme within its own class.

\subsection{Dimensionality Reduction \& Clustering}

\textit{PCA.} PC1 carries 32.1\% of the variance, and its three largest loadings are all operational, which are days\_since\_last\_review (+0.563), amenity\_count (−0.504) and availability\_365 (−0.447). PC2 carries 20.45\% and compares dist\_cbd (+0.786) against host\_listings\_count (−0.556), which separates distance from the CBD from the size of the host's portfolio. PC3 carries 16.8\% and compares accommodates (+0.671) against availability\_365 (−0.591). PC2 explains 20.45\% of the variance and describes geography and scale, but it is unrelated to superhost status. Rather, the variation that tracks superhost status sits almost entirely on PC1.

\textit{K-Means.} K3 is the dormant group, 22.1\% of listings, with a median of 1,883 days since the last review, zero availability across the year, one listing and 23 amenities, and a superhost rate of 6.8\%. K2 is the opposite, at 40 days since the last review and 47 amenities, with the highest superhost rate of the five, 53.9\%. The remaining three sit nearer the 35.72\% base rate for reviewed listings: K1 at 35.3\%, K4 at 43.7\%, and K5 at 40.7\%. They differ mainly in location and size. K1 describes large hosts, a median of 50 listings 1.4 km from the CBD. K4 describes larger listings at 20.6 km. K5 describes small listings at 17.8 km. Here, large and small listings is derived from accommodates, the number of guests that the listing can house.

\textit{Connection to the RQ.} PCA shows that the axis that explains the most is operational, which agrees with the supervised result that fundamentals alone reach 0.166 CV F1 against 0.548 for operational features.

\textit{K-Means versus agglomerative clustering.} Comparing K-Means with agglomerative clustering shows the same structure with different boundaries. The majority overlap is 68.4\%, the share of listings that fall into the largest Ward counterpart of their own K-Means cluster. The Ward profiles carry the same interpretation, each given as size, superhost rate and two distinguishing features. W1 (8,256 listings, 49.9\%, 47 amenities, 56 days since review), W2 (3,829, 30.1\%, 22 amenities, 267 available days), W3 (4,716, 37.5\%, 45 host listings, 1.79 km from the CBD), W4 (3,249, 2.9\%, 2,364 days since review, zero available days) and W5 (1,201, 37.8\%, accommodates 10, 17.55 km from the CBD).

\section{Limitations and Improvement Opportunities}

\textit{Preprocessing.} Price imputation assigns the training-set median of \$245.00 to 6,553 missing rows, and 93.35\% of those are unavailable listings (availability\_365 = 0). Separately, amenity\_count treats every amenity as equally important, so a Wi-Fi router and a pool contribute the same amount. The imputed days\_since\_last\_review value of 4,594 becomes an outlier once standardised for k-NN; modelling that group through has\_review alone, rather than filling a numeric value, was not tried.

An assumption we made is that an availability\_365 of 0 is treated as an unavailable listing, but it could also indicate a fully booked calendar or dates the host has blocked rather than a delisted property.

\textit{Correlation.} Observations are not independent. The 21,251 reviewed listings come from 11,474 hosts, and a host's status is the same across all of their listings.

MI also depends on the number of bins. For days\_since\_last\_review with availability\_365, MI rises from 0.1695 at five bins to 0.3605 at twenty, an increase of 113\%, while NMI moves only from 0.0799 to 0.0896, or 12\%. Comparisons between pairs therefore use NMI.

\textit{Supervised Learning.} The classes are imbalanced. Recall on the superhost class is the weak point, at 0.664 for the tree and 0.560 for k-NN.

Only max\_depth was tuned. For decision trees, max\_depth is set to 10 and min\_samples\_leaf is left at its default of 1, so a leaf can hold a single listing. Nothing other than depth constrains how finely the tree splits, and jointly tuning min\_samples\_leaf was not attempted.

The two feature groups are different sizes. Operational features number eight against three fundamentals, so part of the gap between them may reflect the larger set holding more information rather than that family being more informative.

\textit{Feature Selection.} Continuous features are favoured. A continuous feature offers more possible split points, so it has more chances to reduce impurity. For instance, dist\_cbd ranks 5th by embedded importance at 0.0278 while the filter places it 8th at MI 0.0490. Third place is also very close. availability\_365 and amenity\_count differ by only 0.0015 on mutual information (0.1095 versus 0.1080)

\textit{PCA \& Clustering.} Much of the variance is unexplained. The first three components carry 69.4\%, leaving 30.6\% that none of the interpretations above covers, and PC1 alone accounts for 32.1\%.

The result depends on scaling. The log1p transformation was required because without it long-tailed variables such as host\_listings\_count (maximum 667) and days\_since\_last\_review (maximum 4,593) would dominate the leading components. Different scaling methods may produce different results that are not explored within this report.

The hyperparameter $k$ was also not selected by a formal test. $k=5$ is read from the bend in the elbow curve. SSE falls for every extra cluster, so there is no minimum to look for, and the values for $k=4$ and $k=5$ are 72,991.5 and 64,211.4, a difference of about 12\%. A direct comparison of the $k=4$ and $k=5$ cluster was not carried out.

\FloatBarrier

\section{Conclusion}

The status is partly inferable from the listing record, and the inferable part is operational rather than fundamental. Both feature-selection methods rank days\_since\_last\_review and host\_listings\_count first and second, and the subset experiments supports the conclusion. A decision tree on fundamentals alone reaches 0.166 CV F1 against 0.548 for operational features, while on the test set the values are 0.196 and 0.639. PCA and clustering also support this conclusion, with PC1 accounting for 32.1\% of the variance and K-Means isolating 4,689 low-activity listings at a 6.8\% superhost rate. However, the status can be inferred only in part. The best test F1 score is 0.632 and recall 0.664, finding 1,047 of 1,577 superhost listings while missing 530, and 118 of the 124 superhost listings inactive for over a year remain undetected.

%% GenAI statement inside the 10 pages; References (excluded from the limit) start on page 11
\section*{Generative AI statement}
Throughout this project, we utilized generative AI tools. In coding, it was employed to help consolidate multiple notebooks into the final submission version and review code. For the report, all report text was authored by team members. LLMs was used solely to cross-check numerical consistency across sections and identify errors. All experimental results were executed by the team members locally, and the team assumes full responsibility for the entire submission.

\clearpage
\section*{References}
\begingroup\setlength{\parindent}{0pt}\setlength{\parskip}{3pt}
\hangindent=1.5em\hangafter=1 Inside Airbnb. (2026). Melbourne, Victoria, Australia [Data set]. Retrieved October 2, 2026, from http://insideairbnb.com/get-the-data/

\hangindent=1.5em\hangafter=1 Airbnb. (n.d.). How to become a Superhost. Airbnb Help Centre. Retrieved October 2, 2026, from https://www.airbnb.com.au/help/article/829

\hangindent=1.5em\hangafter=1 Tate, R. F. (1954). Correlation between a discrete and a continuous variable. Point-biserial correlation. The Annals of Mathematical Statistics, 25(3), 603-607. Retrieved October 4, 2026, from https://doi.org/10.1214/aoms/1177728730
\endgroup

\end{document}
